% Совместимость с ninecolors 2022D на Ubuntu 24.04.
%
% tabularray (через doxygen-сгенерированный tim.sty) подтягивает
% ninecolors. ninecolors зовёт \definecolor с моделью «rgb:Hsb», что в
% xcolor 3.01 больше не поддерживается — сборка падает. Сами имена
% (red1..red9 и т.п.) из ninecolors tabularray не использует, поэтому
% безопаснее всего пометить пакет как уже подгруженный — настоящий .sty
% просто не запустится.
\makeatletter
\@namedef{ver@ninecolors.sty}{2022/02/13 v2022D ninecolors stub}
\expandafter\let\csname opt@ninecolors.sty\endcsname\@empty
\makeatother


% Гиперссылки и PDF-метаданные для XeLaTeX.

\usepackage[xetex,
            hyperindex,
            pdffitwindow      = true,
            pdfstartview      = Fit,
            colorlinks        = true,
            linkcolor         = NavyBlue,
            filecolor         = NavyBlue,
            urlcolor          = NavyBlue,
            citecolor         = DarkGreen,
            bookmarksopen     = true,
            bookmarksopenlevel = 0,
            plainpages        = false,
            pdfpagelabels     = true,
            pdftoolbar        = false,
            unicode           = true
           ]{hyperref}

% Метаданные документа.
\hypersetup{
    pdftitle   = {\timproductname},
    pdfauthor  = {\timproductauthor},
    pdfsubject = {TIM developer reference}
}
